\documentclass[11pt]{article}
\usepackage[a4paper,margin=1in]{geometry}
\usepackage{amsmath,amssymb}
\usepackage{pdflscape}

\begin{document}

% Backup material only.  Do not include in the paper unless explicitly needed.
% This file records the complete-D6 open-field cross-line terms in a compact,
% reproducible form.  The explicit machine-readable Borel integrands are stored
% in d6_chunks/*_borel_integrand.wl.

\section*{Backup: open-field cross-line dimension-six terms}

This note records the open-field cross-line completion of the
\(\langle g_s^3G^3\rangle\) sector.  These terms are not included in the
compact coefficient-function appendix because their numerical effect is
negligible in the accepted Borel window and the expressions are better kept in
machine-readable form.

The two cross-line topologies are
\begin{align}
  \Pi^{ij}_{G^3,\mathrm{cross}} &=
  \Pi^{ij}_{c2b1}+\Pi^{ij}_{c1b2},\\
  \Pi^{ij}_{c2b1} &\equiv
  \Pi^{ij}\!\left[S_c^{(GG,\mathrm{open})}S_b^{(G,\mathrm{open})}\right],
  \qquad
  \Pi^{ij}_{c1b2} \equiv
  \Pi^{ij}\!\left[S_c^{(G,\mathrm{open})}S_b^{(GG,\mathrm{open})}\right].
\end{align}
After the spin--1 projection and direct Borel transform, each term is evaluated as
\begin{equation}
  \Pi^{ij}_{X}(M^2,s_0)
  =
  \int_{\xi_-}^{\xi_+}d\xi\,
  I^{ij}_{X}(\xi,M^2),
  \qquad X\in\{c2b1,c1b2\}.
\end{equation}

The explicit \(I^{ij}_{X}\) are stored in the following files:
\begin{center}
\begin{tabular}{lll}
\hline
Channel & Topology & Machine-readable integrand \\
\hline
\(AA\) & \(c2b1\) & \texttt{d6\_chunks/AA\_c2b1\_fast\_borel\_integrand.wl} \\
\(AA\) & \(c1b2\) & \texttt{d6\_chunks/AA\_c1b2\_fast\_borel\_integrand.wl} \\
\(AB\) & \(c2b1\) & \texttt{d6\_chunks/AB\_c2b1\_fast\_borel\_integrand.wl} \\
\(AB\) & \(c1b2\) & \texttt{d6\_chunks/AB\_c1b2\_fast\_borel\_integrand.wl} \\
\(BB\) & \(c2b1\) & \texttt{d6\_chunks/BB\_c2b1\_fast\_borel\_integrand.wl} \\
\(BB\) & \(c1b2\) & \texttt{d6\_chunks/BB\_c1b2\_fast\_borel\_integrand.wl} \\
\hline
\end{tabular}
\end{center}

\subsection*{Explicit cross-line coefficient functions}

We factor the cross-line integrands as
\begin{equation}
  I^{ij}_{X}(\xi,M^2)
  =
  \langle g_s^3G^3\rangle\,
  e^{-\bar s(\xi)/M^2}\,
  D^{ij}_{X}(\xi,M^2),
  \qquad X\in\{c2b1,c1b2\}.
\end{equation}
The explicit coefficient functions are written below in terms of numerators
\(Q^{ij}_{X}\).  This is the same information as the compact machine-readable
files, but split across lines for readability.

\begin{landscape}
\begingroup
\small

\begin{align}
D^{AA}_{c2b1} &=
-\frac{Q^{AA}_{c2b1}}{9216\pi^2M^8(\xi-1)^3\xi^3},&
D^{AA}_{c1b2} &=
\frac{Q^{AA}_{c1b2}}{18432\pi^2M^8(\xi-1)^3\xi^3},\nonumber\\
D^{AB}_{c2b1} &=
-\frac{Q^{AB}_{c2b1}}{9216\pi^2M^8(\xi-1)^4\xi^4(m_b+m_c)},&
D^{AB}_{c1b2} &=
\frac{Q^{AB}_{c1b2}}{18432\pi^2M^8(\xi-1)^4\xi^4(m_b+m_c)},\nonumber\\
D^{BB}_{c2b1} &=
-\frac{Q^{BB}_{c2b1}}{9216\pi^2M^8(\xi-1)^4\xi^4(m_b+m_c)^2},&
D^{BB}_{c1b2} &=
\frac{Q^{BB}_{c1b2}}{18432\pi^2M^8(\xi-1)^4\xi^4(m_b+m_c)^2}.
\end{align}

\begin{align}
Q^{AA}_{c2b1}={}&
2m_b^2(\xi-1)\xi^2
\left[6M^4(\xi-1)^2\xi+M^2m_c^2(-6\xi^2+3\xi+8)+3m_c^4(3\xi-1)\right]\nonumber\\
&+3m_bm_c\xi^2
\left[-24M^4(\xi-1)^2+4M^2m_c^2(\xi^2-3\xi+2)+m_c^4\xi\right]\nonumber\\
&-2\xi^3
\left[-6M^6(\xi-1)^3\xi+6M^4m_c^2(\xi-1)^2\xi
+M^2m_c^4(-3\xi^2+3\xi+5)+3m_c^6\xi\right]\nonumber\\
&+6m_b^4(\xi-1)^2\xi
\left[M^2(\xi^2-1)+m_c^2(2-3\xi)\right]\nonumber\\
&-3m_b^3m_c(\xi-1)\xi
\left[4M^2(\xi^2-4\xi+3)+m_c^2(2\xi-1)\right]
+6m_b^6(\xi-1)^4+3m_b^5m_c(\xi-1)^3,
\end{align}

\begin{align}
Q^{AA}_{c1b2}={}&
6M^6\xi^2(\xi^2-2\xi+6)(\xi-1)^3
+6M^4\xi(\xi^2-2\xi+6)(\xi-1)^2
\left[m_b^2(\xi-1)-m_c^2\xi\right]\nonumber\\
&+M^2(\xi-1)\Big[
m_b^4(\xi-1)^2(3\xi^2-9\xi+1)
-24m_b^3m_c(\xi-1)^2\xi\nonumber\\
&\hspace{1.0cm}
+m_b^2m_c^2\xi(-6\xi^3+27\xi^2-22\xi+1)
+24m_bm_c^3\xi^3+3m_c^4(\xi-4)\xi^3\Big]\nonumber\\
&+3\Big[
m_b^6(\xi-1)^4+2m_b^5m_c(\xi-1)^3
-m_b^4m_c^2(\xi-1)^2\xi(3\xi-2)\nonumber\\
&\hspace{1.0cm}
-2m_b^3m_c^3\xi(2\xi^2-3\xi+1)
+m_b^2m_c^4\xi^2(3\xi^2-4\xi+1)
+2m_bm_c^5\xi^3-m_c^6\xi^4\Big],
\end{align}

\begin{align}
Q^{AB}_{c2b1}={}&
m_b^3(\xi-1)^2\xi^2
\left[36M^4(\xi-1)^2+M^2m_c^2(-72\xi^2+123\xi-56)+9m_c^4(1-3\xi)\right]\nonumber\\
&+m_b^2m_c(\xi-1)\xi^3
\left[M^4(39\xi^2-89\xi+50)+M^2m_c^2(-48\xi^2+92\xi-47)+6m_c^4(1-3\xi)\right]\nonumber\\
&-m_bm_c^2(\xi-1)\xi^3
\left[3M^4(3\xi^2-11\xi+8)+M^2m_c^2(-36\xi^2+39\xi-8)-9m_c^4\xi\right]\nonumber\\
&+m_c\xi^4
\left[22M^6(\xi-1)^2+M^4m_c^2(-21\xi^2+62\xi-41)
+M^2m_c^4(24\xi^2-31\xi+10)+6m_c^6\xi\right]\nonumber\\
&+3m_b^5(\xi-1)^3\xi
\left[4M^2(3\xi^2-7\xi+4)+3m_c^2(3\xi-2)\right]\nonumber\\
&+m_b^4m_c(\xi-1)^2\xi^2
\left[M^2(24\xi^2-61\xi+37)+6m_c^2(3\xi-2)\right]
-9m_b^7(\xi-1)^5-6m_b^6m_c(\xi-1)^4\xi,
\end{align}

\begin{align}
Q^{AB}_{c1b2}={}&
-m_b^3(\xi-1)^2\xi
\left[M^4(\xi-1)^2(15\xi+52)+M^2m_c^2\xi(24\xi^2-23\xi-1)+3m_c^4\xi(3\xi-1)\right]\nonumber\\
&-6m_b^2m_c(\xi-1)\xi^2
\left[3M^4(\xi-1)^2(7\xi-3)+M^2m_c^2(8\xi^3+3\xi^2-18\xi+7)+m_c^4\xi(3\xi-1)\right]\nonumber\\
&+m_b(\xi-1)\xi^2
\left[104M^6(\xi-1)^3+4M^4m_c^2(\xi-1)^2(6\xi+13)
+4M^2m_c^4\xi(3\xi^2-\xi-2)+3m_c^6\xi^2\right]\nonumber\\
&+6m_c\xi^3
\left[24M^6(\xi-1)^3+12M^4m_c^2(\xi-1)^2(2\xi+1)
+4M^2m_c^4\xi(\xi^2+\xi-2)+m_c^6\xi^2\right]\nonumber\\
&+m_b^5(\xi-1)^3\xi
\left[M^2(12\xi^2-19\xi+7)+3m_c^2(3\xi-2)\right]\nonumber\\
&+6m_b^4m_c\xi
\left[M^2(4\xi+7)(\xi-1)^4+m_c^2\xi(3\xi-2)(\xi-1)^2\right]
-3m_b^7(\xi-1)^5-6m_b^6m_c(\xi-1)^4\xi,
\end{align}

\begin{align}
Q^{BB}_{c2b1}={}&
2m_b^4(\xi-1)^2\xi^2
\left[6M^4(\xi-6)(\xi-1)^2+M^2m_c^2(54\xi^2-102\xi+53)+9m_c^4(2\xi-1)\right]\nonumber\\
&+m_b^3m_c(\xi-1)\xi^2
\left[2M^4(27\xi^3-85\xi^2+100\xi-42)+2M^2m_c^2(54\xi^2-83\xi+32)+9m_c^4(3\xi-1)\right]\nonumber\\
&+2m_b^2(\xi-1)\xi^3
\left[12M^6(\xi-1)^3\xi-6M^4m_c^2(2\xi^3-10\xi^2+19\xi-11)
-2M^2m_c^4(27\xi^2-39\xi+17)+3m_c^6(1-4\xi)\right]\nonumber\\
&-m_bm_c\xi^3
\left[-4M^6(\xi-12)(\xi-1)^2+2M^4m_c^2(27\xi^3-58\xi^2+55\xi-24)
+2M^2m_c^4(27\xi^2-28\xi+4)+9m_c^6\xi\right]\nonumber\\
&+2\xi^4
\left[12M^8(\xi-1)^4\xi-6M^6m_c^2(\xi-1)^2(2\xi^2-2\xi-5)
+6M^4m_c^4(\xi^3-2\xi^2+6\xi-5)\right.\nonumber\\
&\hspace{1.0cm}\left.
+M^2m_c^6(18\xi^2-18\xi+5)+3m_c^8\xi\right]\nonumber\\
&-6m_b^6(\xi-1)^3\xi
\left[2M^2(3\xi^2-7\xi+4)+m_c^2(4\xi-3)\right]\nonumber\\
&-m_b^5m_c(\xi-1)^2\xi
\left[2M^2(27\xi^2-55\xi+28)+9m_c^2(3\xi-2)\right]
+6m_b^8(\xi-1)^5+9m_b^7m_c(\xi-1)^4,
\end{align}

\begin{align}
Q^{BB}_{c1b2}={}&
12M^8\xi^3(\xi^2-2\xi+6)(\xi-1)^4\nonumber\\
&+2M^6\xi^2(\xi-1)^3
\left[3m_b^2(2\xi^3-6\xi^2-19\xi+23)-8m_bm_c(2\xi+13)
-6m_c^2\xi(\xi^2-2\xi-12)\right]\nonumber\\
&+2M^4\xi(\xi-1)^2
\Big[3m_b^4(\xi-1)^2(\xi^2+4\xi+5)
+4m_b^3m_c(8\xi^2+5\xi-13)\nonumber\\
&\hspace{1.0cm}
+3m_b^2m_c^2\xi(-2\xi^3-12\xi^2+3\xi+11)
-4m_bm_c^3\xi(8\xi+13)+3m_c^4\xi^2(\xi^2+10\xi+6)\Big]\nonumber\\
&-M^2(\xi-1)\left[m_b^2(\xi-1)-m_c^2\xi\right]^2
\left[m_b^2(18\xi^3-30\xi^2+11\xi+1)+16m_bm_c\xi-6m_c^2\xi^2(3\xi+2)\right]\nonumber\\
&+3\left[m_b^2(\xi-1)-m_c^2\xi\right]^3
\left[m_b^2(\xi-1)^2-m_c^2\xi^2\right].
\end{align}

\endgroup
\end{landscape}

At the central point \(M^2=8.0~\mathrm{GeV}^2\), \(s_0=54.0~\mathrm{GeV}^2\),
the cached cross-line moments are
\begin{center}
\begin{tabular}{lrrr}
\hline
Channel & \(\Pi_{c2b1}\) & \(\Pi_{c1b2}\) & Sum \\
\hline
\(AA\) & \(-1.67570\times 10^{-6}\) & \(3.37871\times 10^{-7}\) & \(-1.33783\times 10^{-6}\) \\
\(AB\) & \(6.90753\times 10^{-7}\) & \(1.19214\times 10^{-7}\) & \(8.09966\times 10^{-7}\) \\
\(BB\) & \(-1.16841\times 10^{-6}\) & \(-7.55718\times 10^{-7}\) & \(-1.92413\times 10^{-6}\) \\
\hline
\end{tabular}
\end{center}
Adding these terms changes the central mixing angle by
\begin{equation}
  \theta_{\mathrm{with\,cross}}-\theta_{\mathrm{single-line}}
  =
  -2.84\times10^{-4}\ \mathrm{deg}.
\end{equation}
The corresponding relative shifts of the three Borel moments are
\begin{equation}
  \left(
  \frac{\Pi^{AA}_{\mathrm{cross}}}{\Pi^{AA}_{\mathrm{base}}},
  \frac{\Pi^{AB}_{\mathrm{cross}}}{\Pi^{AB}_{\mathrm{base}}},
  \frac{\Pi^{BB}_{\mathrm{cross}}}{\Pi^{BB}_{\mathrm{base}}}
  \right)
  =
  (-2.84,\,-2.37,\,-4.47)\times10^{-5}.
\end{equation}

\end{document}
